\section{De la Architecture al Universal Substrate}

La conferencia termina situando de nuevo al Transformer dentro de la consolidation thesis. Ya se extendió desde el translation backbone hacia vision, contrastive learning, modelos de lenguaje y tool-using systems, pero el curso también enumera de forma explícita lo que no cubrió: layer norm, modificaciones a la FFN, higher-order attention y faster decoding. Universal substrate describe una tendencia; no es una declaración de que «todos los problemas ya están unificados».

\subsection{Extensiones a otros campos y omisiones deliberadas de esta clase}

El Vision Transformer trata los patches de la imagen como tokens, y el contrastive model alinea representaciones entre modalidades; estos cambios demuestran que el block de token mixing + feature transformation es transferible. Al mismo tiempo, la pre/post layer norm, la gated FFN, la higher-order interaction y el speculative/non-autoregressive decoding influyen de manera notable en los sistemas modernos; esta clase solo sigue la línea principal de attention/position/system.

\lecturefigure{slide-37-advancements-other-areas.jpg}{El Transformer se extiende a vision y contrastive learning}{Diapositiva V037 recuperada de la grabación oficial; 00:52:33}

\begin{knowledgebox}{Lectura de la figura: lo que se transfiere es el block motif, no que las entradas sean iguales por naturaleza}
Las imágenes requieren patch/tokenización y position bidimensional; el contrastive learning requiere un objective entre muestras. Primero observe el Transformer encoder compartido y después el input/output específico de cada modalidad (modality-specific); el éxito entre campos respalda la architecture reuse, pero no implica que deba eliminarse todo inductive bias.
\end{knowledgebox}

\lecturefigure{slide-38-not-covered-topics.jpg}{Direcciones de evolución del Transformer que esta clase no desarrolla, pero que son igual de importantes}{Diapositiva V038 recuperada de la grabación oficial; 00:54:15}

\begin{warningbox}{Una conferencia de 80 minutos no sustituye un mapa completo del campo}
La pre-vs-post norm afecta la estabilidad de las redes profundas, la FFN/gating afecta la capacity, la higher-order attention cambia la interaction y el speculative decoding mejora la inference. Interpretar «no cubierto» como «no importante» haría que los apuntes fueran más tajantes que la propia clase.
\end{warningbox}

\teachervoice{Que el ponente muestre por iniciativa propia la diapositiva Not Covered es una forma de evidence hygiene: reconoce que muchos avances importantes no entraron en el relato. Unos apuntes de calidad también deben conservar la delimitación del alcance, en lugar de meter a la fuerza en la clase todos los términos posteriores solo para parecer completos.}

\subsection{El LLM como Prompt-programmable Substrate}

Cuando un mismo pretrained model puede ejecutar una gran variedad de tareas mediante el prompt, empieza a parecerse a un substrate de cómputo general: el usuario ya no entrena un modelo independiente para cada tarea, sino que especifica el comportamiento mediante el context. Aquí la «programming» sigue siendo probabilistic conditioning, limitada por los datos de entrenamiento, el instruction following, la context window y la reliability.

\lecturefigure{slide-39-llms-universal-substrate.jpg}{El LLM como prompt-programmable universal substrate}{Diapositiva V039 recuperada de la grabación oficial; 00:57:15}

\begin{knowledgebox}{Lectura de la figura: el prompt es una interfaz, no una prueba de confiabilidad}
El diagrama al estilo Pac-Man expresa que un solo modelo absorbe múltiples tareas. Primero pregunte si las tareas realmente comparten los model weights y después si la salida puede verificarse; la few-shot breadth muestra consolidation, pero no garantiza exact computation, fresh knowledge, long-horizon execution ni safety.
\end{knowledgebox}

\begin{importantbox}{La consolidation traslada los problemas hacia etapas posteriores}
Cuando la interfaz del modelo se unifica, las diferencias no desaparecen: se trasladan a la data mixture, el prompt/tool schema, la evaluation, el serving, el feedback y la safety. Lo que se simplifica es la API superficial, no necesariamente la pila de ingeniería completa.
\end{importantbox}

\subsection{El Transformer moderno es un Data-center System}

El entrenamiento y el serving de los modelos grandes abarcan miles de GPU e implican \term{collectives (comunicación colectiva)} como all-reduce, all-gather y reduce-scatter, además de la topología de red, la congestión, las estrategias de paralelismo y la tolerancia a fallas. Un layer en el diagrama del modelo, en tiempo de ejecución, se divide en parámetros, activations y un communication schedule repartidos entre varios devices.

\lecturefigure{slide-40-modern-transformer-datacenter.jpg}{El Transformer moderno a gran escala debe entenderse como un data-center system}{Diapositiva V040 recuperada de la grabación oficial; 00:59:06}

\begin{knowledgebox}{Lectura de la figura: al crecer el número de parámetros, la red también entra en el critical path del modelo}
La fotografía del centro de datos no es decorativa. Primero pregunte cómo se hace la partición con tensor/data/pipeline parallel y después si los collectives se superponen con el compute (compute overlap); la quality del modelo puede ser la misma, pero la cluster topology y la congestion determinan si el entrenamiento es escalable. Un architecture researcher también necesita entender los systems boundaries.
\end{knowledgebox}

\teachervoice{El ponente dice que, si de verdad se investiga el modern large-scale Transformer, hay que estudiar all-reduce, InfiniBand, RoCE y congestion. El Transformer grande «now in some sense is a data center»; esto cierra el ciclo con la inversión del juicio sobre la machine capacity de Dartmouth planteado al inicio.}

\subsection{Tools: qué capacidades deben estar dentro del modelo y cuáles deben recurrir a sistemas externos}

El \term{tool use (uso de herramientas)} permite que el modelo invoque, mediante una interface estructurada, una calculadora, búsqueda, una base de datos, un code executor u otro service. No es un regreso al pipeline frágil, sino dejar fuera las capacidades de software verificables, exactas y ya existentes, mientras el modelo se encarga de elegir, combinar e interpretar.

\lecturefigure{slide-41-llms-using-tools.jpg}{El LLM se conecta a búsqueda, cálculo y software externo mediante herramientas}{Diapositiva V041 recuperada de la grabación oficial; 00:59:24}

\begin{knowledgebox}{Lectura de la figura: la Swiss-army knife sigue necesitando interfaces claras}
Cada herramienta tiene su schema, sus permisos, su cost y sus failure modes. Primero observe cuándo decide el modelo invocarla y después cómo se verifica el resultado y cómo pasa al siguiente paso; la tool breadth puede ampliar la capability, pero un routing erróneo, la prompt injection y las actions irreversibles también amplían el riesgo.
\end{knowledgebox}

\teachervoice{En la sesión de Q\&A, Vaswani ofrece un juicio muy ingenieril: no conviene gastar «mil millones de FLOPs por posición» para calcular dos números; ciertas capacidades deben delegarse a herramientas externas, y el modelo debe encargarse del thinking que requiere juicio semántico y combinación. El límite de las herramientas es una capability allocation, no una señal de que el modelo sea poco puro.}

\subsection{It Is Time to Build: la retroalimentación del producto también es una señal de aprendizaje}

El ponente sostiene que la interacción con usuarios reales puede exponer carencias de capacidad y formar un ciclo de data/feedback/model improvement. El producto no genera buenos datos de forma automática: la retroalimentación tiene selection bias, problemas de privacy, reward hacking y KPI de corto plazo; pero sin workflows reales, es difícil que los investigadores sepan en qué interfaces falla realmente el modelo.

\lecturefigure{slide-42-time-to-build.jpg}{El uso real y la mejora del modelo pueden formar un ciclo cerrado de retroalimentación del producto}{Diapositiva V042 recuperada de la grabación oficial; 01:00:06}

\begin{importantbox}{El valor de investigación de Build proviene de las Failures observables}
Un ciclo cerrado de alto valor debe registrar la tarea, la trayectoria de herramientas, las correcciones del usuario, el resultado de la ejecución y la clasificación de fallas, y distinguir entre «al usuario le gusta» y «la tarea es correcta». Solo al convertir la retroalimentación en artifacts de entrenamiento/evaluation auditables, el uso del producto mejorará el modelo, en lugar de solo amplificar los sesgos existentes.
\end{importantbox}

\subsection{Resumen de la sección}

La consolidation del Transformer ya atraviesa modalidades y tareas, pero todavía queda mucho trabajo de architecture sin cubrir. Una vez que el LLM se convierte en un substrate general, los problemas de sistemas se desplazan hacia el data center, las tools, el feedback y la verification; el modelo unificado no hizo desaparecer la ingeniería, sino que reorganizó sus límites.
